% Working draft, not a report of completed runs. The Overleaf project was not
% accessible when this file was prepared. Resolve the execution notes in
% README.md and insert measured results before integrating into the paper.
\section{Experiments}
\label{sec:experiments}

\paragraph{Synthetic corpus.}
We generated 12,000 clarinet score--performance pairs: 10,000 from
procedurally generated scores and 2,000 from snippets of three existing
scores. Each pair contains a clean MusicXML score and a performance with
one to eight injected content errors drawn from wrong notes, missed notes,
extra notes, rhythm errors, and intonation errors, with optional repetitions.
Reference and faulty audio are rendered with the FreePats clarinet SoundFont
at 22,050\,Hz in mono. The faulty recordings total 134.08 hours. Error
annotations, generation-time note identities, MIDI event timing, audio
alignments, and mel spectrograms are saved with each pair.
We retain 11,165 pairs after verifying synthetic note lineage, audible tie
groups, and agreement between MIDI-derived labels and the renderer's event
times. We exclude 835 pairs with ambiguous supervision or inconsistent
rendering times. The retained corpus is divided into 10,049 training and
1,116 validation pairs, with the same membership for each compared method.
% This split is frozen for the ongoing baseline runs. ALIGN must use it too.

\paragraph{Evaluation setting.}
We compare ALIGN, Polytune, and LadderSym on a common validation set of
score--performance pairs. Each example consists of a clean symbolic score,
a clarinet performance recording, and error annotations. Split membership
is fixed across methods. Training examples are used for parameter updates,
and validation examples are used for checkpoint selection and calibration.
The reported scores therefore describe validation performance under this
selection procedure. All methods are evaluated on the same complete list
of validation examples.
% Insert audited dataset identity and actual training/validation counts for
% the new corpus. Do not reuse the legacy align_v1 split or exclusion counts.
% If using a per-clip split, retain the following qualification:
The split is defined over clips; these results characterize performance
within the sampled score collection and do not establish generalization
to unseen pieces.

\paragraph{Baselines.}
Polytune~\cite{chou2025polytune} compares performance audio with a synthesized
audio rendering of the reference score using a Transformer encoder--decoder.
It produces note events labeled as correct, missing, or extra. A wrong pitch
is represented by a missing expected note and an extra performed note.
LadderSym~\cite{chou2026laddersym} introduces cross-attention between the
reference and performance encoder streams and supplies symbolic score
tokens as a decoder prompt. We use the prompted LadderSym variant, which
has access to both the score rendering and reference MIDI. Both baselines
use the authors' implementations with the ALIGN data loaders and
class-preserving MIDI export. Their reference audio and MIDI retain the
score timeline, and performance and reference windows are selected at the
same absolute time. This input convention can limit their tolerance to
large tempo changes or repeated passages.

\paragraph{ALIGN implementation.}
We evaluate the transcription-first implementation. A frozen Basic Pitch
0.4.0 front end transcribes the performance into note events, with pitch
conversion to the written-score convention. The decoder applies clarinet
range and harmonic filtering and merges weak-onset splits of sustained
notes. Layer~1 detects repeated phrases in the transcribed sequence using
tempo-tolerant matching. A learned multilayer perceptron scores repetition
candidates and provides a one-note recovery rule when no longer repetition
has been found. Layer~2 aligns the first-pass performance notes to the clean
score with a contextual note aligner; replayed notes reuse the source
phrase's score correspondence. The aligner uses separate two-layer
bidirectional GRUs with 64 hidden units per direction, contextual pair
scores, and an extra-note output. A monotonic sequence decoder permits
matched or substituted notes, insertions, and deletions. Tied score notes
are collapsed before alignment.

Layer~3 checks note durations using these same repetition-aware
correspondences. For each usable note, the performed-to-score duration
ratio is compared with the median ratio for the clip. The default detector
requires an absolute log-ratio deviation of at least 0.35 and a duration
deviation of at least 120\,ms, and it requires at least four usable pairs.
The output comprises wrong notes, extra notes, missed notes, repetitions,
and rhythm errors, localized on the clean score. Intonation detection and
the optional timbre stage are disabled in this configuration. The
performance score and synthetic lineage maps provide training supervision;
they are not inputs to inference.

\paragraph{Training and selection.}
% Proposed run settings below must be reconciled with actual run artifacts.
Both baselines are initialized from scratch. Their training configuration
uses seed 365, weighted token cross-entropy with error-loss coefficient 8,
and a learning-rate schedule configured for 40 epochs. Training was stopped
early; the selected checkpoints correspond to 31 epochs for Polytune and
26 epochs for prompted LadderSym. The learning rate has a peak of $2\times10^{-5}$, a
4,000-update warmup, and cosine decay to half the peak. The decoder event
budget is 1,024 tokens, and LadderSym additionally uses a 1,024-token
symbolic prompt budget. Checkpoints are selected by validation loss.
Validation loss uses deterministically selected segments; the F1 evaluation
covers each complete validation recording.
Each batch contains one clip contributing one sampled 2.048-second segment.
Both models use 16 gradient-accumulation steps and therefore accumulate
16 clips per optimizer update, except for a final incomplete
accumulation. Training uses mixed bfloat16 precision on RTX A5500 GPUs.
Inference uses float32 greedy decoding. For prompted LadderSym, generation
appends the same unmasked decoder start token after the padded score prompt
as training. Decoder key--value caching is enabled after verifying cached
and uncached generation consistency.
% Record actual initialization, effective batch sizes, accumulation, completed
% epochs, hardware, and selected checkpoint hashes after training.

ALIGN's repetition scorer uses AdamW with learning rate $8\times10^{-4}$
and weight decay $10^{-3}$. Its checkpoint and detection threshold are
selected on validation data. The contextual aligner is first trained on
exact synthetic note correspondences, removing replay copies from the
first-pass sequence, using learning rate $5\times10^{-4}$. It is then
fine-tuned on cached Basic Pitch outputs after Layer~1 using learning rate
$2\times10^{-4}$. Fine-tuning combines note-assignment cross-entropy with
a monotonic sequence-level alignment loss. Both alignment stages use
AdamW with weight decay $10^{-3}$ and select checkpoints by validation
note-assignment accuracy. Basic Pitch remains frozen throughout.
% Insert actual training sample counts, epochs, batches, calibrated decoder
% settings, and the validation decision to retain a candidate aligner.

\paragraph{Metrics.}
The two baselines natively predict note classes, whereas ALIGN predicts
typed error regions. We therefore distinguish note-level error detection
from error-region localization. For the baselines' native output, we
report precision, recall, and F1 separately for extra, missing, and correct
notes using one-to-one onset-and-pitch matching with 50\,ms onset tolerance
and 50-cent pitch tolerance. Note offsets are not required to match.
Class identity is read from MIDI track names, including when a class is
empty. Micro scores pool matched, predicted, and reference note counts;
per-clip mean scores are reported separately. Pooling all three classes
without their identities measures transcription and is not a substitute
for error-class F1.
We also report an overall class-aware micro-F1 by pooling the three
class-specific match counts. Notes emitted without a valid class receive
no class-match credit and count as false positives in this overall score.

For a common error-region comparison, baseline note predictions are
converted using the same fixed adapter for both models. Mutually nearest
extra and missing onsets within 100\,ms form a wrong-note prediction;
unpaired notes yield extra-note or missed-note predictions. Repetition
regions are inferred from runs of extra notes whose pitch sequences match
preceding performed material. These repetition outputs are produced by
the adapter, rather than a native repetition class. The shared comparison
covers wrong notes, extra notes, missed notes, and repetitions, with
one-to-one onset matching evaluated separately by error type. Primary
onset tolerance is 50\,ms; 100 and 200\,ms provide tolerance sensitivity.
Ground truth retains first-pass errors and the repetition annotation,
excluding duplicated error annotations on replayed passes. Rhythm is
reported separately for ALIGN because the baseline output vocabularies
do not contain a rhythm-error class.
% This common adapter/evaluator must be validated on the actual annotation
% schema before filling a comparison table. In particular, schema 1.2
% score-melody windows must not be assumed to be wall-clock error onsets.

ALIGN's score-localization diagnostic uses its implemented melody-set F1.
Predicted and reference pitch sequences are matched one-to-one. A range
match requires an equal sequence or an LCS-Dice similarity of at least
0.80 together with a length ratio of at least 0.60; containment alone is
insufficient. A range match receives weight 1 for the correct error type
and 0.5 for a different type. Clips with empty prediction and reference
sets receive a score of 1. This score is reported separately from the
onset-based metrics above.

\paragraph{Validation results.}
% Native baseline predictions are complete; the common ALIGN comparison is pending.
Table~\ref{tab:validation-native-baselines} reports the native note-level
results for both baselines on all 1,116 validation pairs. The final comparison with ALIGN requires
completed predictions and a verified common error-region evaluation.

\begin{table}[t]
\centering
\caption{Native note-level micro-F1 (\%) on the common validation set.
Overall retains class identity and counts unclassified notes as false positives.}
\label{tab:validation-native-baselines}
\begin{tabular}{lcccc}
\hline
Method & Extra & Missing & Correct & Overall \\
\hline
Polytune & 17.15 & 0.94 & 18.41 & 16.99 \\
LadderSym (prompted) & 7.79 & 2.26 & 16.15 & 11.66 \\
\hline
\end{tabular}
\end{table}
